\documentclass{article}

\usepackage{preamble}

\title{Advanced Radio License Notes}
\author{Oskar Meszaros (z5636164)}
\date{July 2026}

\begin{document}

\maketitle

\tableofcontents

\section{Introduction}
This are the notes from following the advanced course on HamRadio.au.

\section{Electronics and Circuit Theory}
\subsection{Impedance}
\subsubsection{Introduction}
Impedance is how a circuit stores or resists AC (or any changing) energy. If the impedances within a circuit aren't matched, then power gets reflected back instead of going where it is needed.

Impedance is the combination of resistance and reactance. Resistance is the restriction of the "flow" of energy, and converts that restricted energy to heat. This is measured in Ohms $\Omega $. Reactance can be thought of as a flywheel, it stores energy temporarily and releases it. This is what capacitors and inductors do. This is also measured in Ohms. This means that impedance is the total opposition to current.
\begin{callout}{info}[Why this matters]
The transmitter expects to see $50\Omega $ of impedance. The coax is also $50\Omega $. If the antenna isn't also close to $50\Omega $, then the power from the transmitter gets reflected, which is what SWR (standard wave ratio) measures.
\end{callout}

\subsubsection{Reactance}
There are two types of reactance. One is found in inductors, and the other in capacitors.

Inductors (any coil of wire) opposes changes in current, i.e. their reactance increases with frequency
\[
    X_L=2\pi fL.
\]
This means that the higher the frequency, the more and inductor blocks the signal. This is why they are used in low-pass filters, they all for low frequencies but block high ones.

Capacitors opposes changes in voltage. Their reactance decreases with frequency
\[
    X_C=\frac{1}{2\pi fC}.
\]
The higher the frequency, the easier it passes through the capacitors. This is why a capacitor in series blocks DC but passes RF.

\subsubsection{Impedance formula}
We can then write impedance as
\[
    Z=R+jX,
\]
where $R$ is the resistance, $X$ is the impedance, and $j$ is used in place of $i$ for imaginary numbers due to $i$ being used for current. The resistance part is where the energy is consumed, and the reactive part is where the energy is stored and returned.
The total impedance is then given by the magnitude of this complex number
\[
    |Z|=\sqrt{R^2+X^2}
\]

\begin{callout}{Example}[Example]
    An antenna has $R=50\Omega $ and $X=25\Omega $. The total impedance magnitude is
    \begin{align*}
        Z &= \sqrt{2500+625} \\
          &=\sqrt{3125} \\
          &\approx 56\Omega 
    \end{align*}
\end{callout}

\subsection{Resonance}
\subsubsection{Introduction}
When an inductor and capacitor are combined, at one specific frequency their reactances become equal and therefore cancel each other out. This is called resonance, and when this happens the circuit behaves as purely resistive. This is how a single frequency is selected while rejecting everything else.
\begin{callout}{Example}[Radio Example]
    Tuned circuits within a receiver use this resonance to select a station to hear, and reject everything else. Antennas are resonant at a specific frequency, which is where it works best.
\end{callout}

The formula to find the resonant frequency is given by
\[
    f_r=\frac{1}{2\pi \sqrt{LC}}.
\]
This says that for a given inductor $L$ and capacitor $C$, there's exactly one frequency where they resonate. This is an important formula.

\subsubsection{Series and Parallel Resonance}
Series and parallel arrangements of $L$ and $C$ result in opposite behaviour at resonance.
\begin{table}[h]
    \centering
    \caption{Behaviour of circuits at resonance}
    \label{tab:1}
    \begin{tabular}{c|c|c}
        & Series $LC$ & Parallel $LC$ \\
        \hline
        \hline
        Impedance & Minimum (just $R$) & Maximum \\
        Current & Maximum & Minimum from source \\
        Used for & Passing a desired frequency (band-pass) & Blocking a desired frequency (or selecting it) (band-stop)
    \end{tabular}
\end{table}

\begin{callout}{Example}[Radio Example]
    A series-resonant circuit in an antenna trap passed the trap's frequency through, while a parallel-resonant circuit in a receiver's IF stage has high-impedance at the desired frequency, which develops a large signal voltage across it.
\end{callout}

\subsubsection{Q-Factor}
The Q-factor described how selective a resonant circuit is. A higher A means a very sharp tuning (with narrow bandwidth), and low Q means a very broad tuning (wide bandwidth)
\[
    Q=\frac{f_r}{\text{Bandwidth}}
\]
For a series $RLC$, $Q=\frac{X_L}{R}$, so less resistance means higher Q.
\begin{callout}{info}[Bandwidth Rule]
    Bandwidth is $\frac{f_r}{Q}$, so a circuit with $Q=100$ at 7 MHz has a bandwidth of 70 kHz, which is enough to cover a portion of the 40m band. A crystal filter width $Q=50 000$ at 9 MHz has a bandwidth of only 180 Hz, which is perfect for CW.
\end{callout}



\end{document}
